% Flavor catalog per-process entry.
% process_id: CR005
% family: collider_rs
% owner_agent_id: PKA-CR005
% writer_agent_id: WA
% checker_agent_id: null
% opus_signoff_id: null
% Canonical metadata sidecar: CR005.yaml
% Status is controlled by the YAML sidecar status_history, not this comment.

\section{CR005: \(pp \to (\gamma^{(1)},Z^{(1)})_{\rm KK}\to \ell^+\ell^-\)}

\subsection*{Process}
\textbf{Standard notation:}
\[
  pp \to (\gamma^{(1)},Z^{(1)})_{\rm KK}\to \ell^+\ell^-,
  \qquad \ell=e,\mu .
\]
This entry covers prompt high-mass dielectron and dimuon resonance searches
used as proxies for neutral electroweak KK gauge excitations.  The
experimental papers usually phrase the search in terms of spin-1 \(Z'\)
benchmarks such as \(Z'_{\rm SSM}\), not as a complete custodial-RS model.

\subsection*{Current best limit(s)}
The canonical current benchmark is the CMS full-Run-2 high-mass dilepton
search, which excludes a sequential-standard-model spin-1 \(Z'_{\rm SSM}\)
below
\[
  m_{Z'_{\rm SSM}} = 5.15~{\rm TeV}
  \quad (95\%~{\rm CL},~ee+\mu\mu).
\]
CMS also reports \(m_{Z'_\psi}>4.56~{\rm TeV}\) for the corresponding
superstring-inspired narrow benchmark.  The closest ATLAS full-Run-2 result
sets \(m_{Z'_{\rm SSM}}>5.1~{\rm TeV}\) in the combined dilepton channel.
ATLAS further gives model-independent fiducial
\(\sigma_{\rm fid}\times{\cal B}\) upper limits that reach about
\(0.014~{\rm fb}\) at a 6 TeV pole mass for a zero-width signal.  The PDG
\(Z'\)-searches review summarizes the ATLAS/CMS prompt dilepton searches as
setting 95\% CL upper cross-section limits as low as \(0.02~{\rm fb}\).

These limits are sidecar-measured observables, not theory predictions; see
\texttt{flavor\_catalog/references/CR005/cms\_2021\_arxiv2103\_02708.txt},
\texttt{atlas\_2019\_arxiv1903\_06248.txt}, and
\texttt{pdg2024\_zprime\_searches.txt}.

\subsection*{Relevance to RS}
For an RS model with bulk electroweak gauge fields, this search constrains
the first neutral electroweak KK tower, schematically
\((\gamma^{(1)},Z^{(1)},Z'_{\rm custodial})\), only after specifying the
light-quark production coupling, the lepton coupling, the total width, and
the branching fraction to \(e^+e^-\) and \(\mu^+\mu^-\).  A sequential
\(Z'\) limit is therefore a useful collider benchmark, but it is not a
model-independent lower bound on the RS \(M_{\rm KK}\).

The distinction matters for the anarchic-flavor scan.  The current quark
scan methodology note quotes the low-energy \(\Delta F=2\) headline crossing
\(M_{\rm KK}^{\min}=47.26~{\rm TeV}\) at \(g_*=3\) and 50\% acceptance.
The direct dilepton reach is an order of magnitude lower in the SSM-like
benchmark and can weaken further in custodial RS if the neutral KK bosons
couple weakly to light quarks or leptons, are broad, or preferentially decay
to \(t\bar t\), \(b\bar b\), \(W_LW_L\), \(Zh\), Higgs-sector states, or
fermion partners.  CR005 is therefore a PRIMARY collider cross-check and a
handle on non-anarchic or nonstandard coupling patterns, not the leading
constraint for the anarchic-flavor RS baseline.

\subsection*{Post-2010 developments}
The post-2010 search history is a steady luminosity and energy progression.
ATLAS \texttt{arXiv:1209.2535} is the Run-1 reference most directly naming
Kaluza--Klein \(Z/\gamma\) bosons among the interpreted high-mass dilepton
signals.  ATLAS \texttt{arXiv:1707.02424} and CMS
\texttt{arXiv:1803.06292} moved the same prompt \(ee/\mu\mu\) strategy to
the early 13 TeV Run-2 dataset and set the predecessor high-mass resonance
limits.  ATLAS \texttt{arXiv:1903.06248} then used the full Run-2 ATLAS
sample and provided generic width-dependent fiducial cross-section limits
plus \(Z'_{\rm SSM}\), \(Z'_\chi\), \(Z'_\psi\), and HVT interpretations.
CMS \texttt{arXiv:2103.02708} used the full Run-2 CMS electron and muon
datasets, supplied the sharpest combined \(Z'_{\rm SSM}\) mass benchmark for
this entry, and also published recast-friendly coupling-plane information in
HEPData.

\subsection*{Validity / model dependence}
The measured collider objects are cross-section upper limits and benchmark
mass exclusions.  Translating them into custodial-RS parameters requires
model input: production is usually \(q\bar q\)-initiated, the \(Z'\) or KK
state is treated as a narrow or moderately broad spin-1 resonance, and the
quoted SSM benchmark assumes SM-like fermion couplings.  In realistic
custodial RS, the first neutral electroweak resonances can have suppressed
light-fermion couplings and enhanced couplings to IR-localized third
generation fermions and longitudinal electroweak bosons.  The dilepton
branching fraction may therefore be much smaller than in \(Z'_{\rm SSM}\).
Conversely, a non-anarchic flavor structure or a model with deliberately
large lepton couplings can make this channel more competitive than the
minimal custodial-RS expectation.

\subsection*{Code coverage}
\textbf{NO.}  Required greps over \texttt{quarkConstraints/}, \texttt{qcd/},
\texttt{flavorConstraints/}, \texttt{neutrinos/}, \texttt{yukawa/},
\texttt{warpConfig/}, \texttt{solvers/}, \texttt{scanParams/}, and
\texttt{tests/} found no direct LHC dilepton-resonance likelihood, no
ATLAS/CMS/HEPData collider-limit ingestion, and no neutral electroweak KK
filter.  The adjacent code is low-energy or spectrum bookkeeping:
\texttt{quarkConstraints/deltaf2.py:1--10} describes a tree-level
KK-gluon-inspired \(\Delta F=2\) benchmark, and
\texttt{quarkConstraints/deltaf2.py:320--324} matches a compact meson-mixing
operator basis at \(M_{\rm KK}\).  \texttt{quarkConstraints/couplings.py:96--124}
builds quark mass-basis KK-gluon couplings, not electroweak dilepton
production.  The only collider-looking grep hit was
\texttt{tests/test\_alpha\_s.py:88--90}, a CMS/RunDec \(\alpha_s\) example,
not an LHC search constraint.

\subsection*{Implementation difficulty}
\textbf{HIGH.}  A faithful implementation is not a one-line mass cut.  It
needs a signal model for the neutral electroweak KK spectrum and couplings,
parton luminosities or event generation, widths and branching fractions, and
a likelihood or external reinterpretation of the ATLAS/CMS dilepton spectra
and limits.  A practical implementation would likely use exported HEPData
tables plus a recasting tool or a validated internal likelihood, and would
need separate branches for narrow \(Z'\)-like, broad, and third-generation
dominated custodial-RS scenarios.

\subsection*{Key references}
Process-local source keys before bibliography consolidation are
\texttt{CMS2021\_HighMassDilepton},
\texttt{ATLAS2019\_HighMassDilepton},
\texttt{PDG2024\_ZPrimeSearches},
\texttt{HEPDataCMS2021\_HighMassDilepton},
\texttt{HEPDataATLAS2019\_HighMassDilepton},
\texttt{DavoudiaslHewettRizzo1999\_BulkGauge},
\texttt{AgasheDelgadoMaySundrum2003\_CustodialRS},
\texttt{AgasheServant2004\_WarpedGUT},
\texttt{BellaEtAl2010\_KKEWGauge},
\texttt{ATLAS2012\_Run1Dilepton},
\texttt{ATLAS2017\_Run2Dilepton36fb}, and
\texttt{CMS2018\_Run2Dilepton36fb}.
